\documentclass[report, twoside, UTF8, AutoFakeBold = 1, AutoFakeSlant, zihao = -4]{config}

% 封面图片定义
\def \titlePageImages{
    \includegraphics[width=0.45\textwidth] {NWPU-logo.pdf}\\ % 校徽
    \includegraphics[width=0.45\textwidth] {NWPU-title-CN.pdf}\\ % 中文校名
    \includegraphics[width=0.55\textwidth] {NWPU-title-EN.pdf}\\ % 英文校名    
}

% 文档信息定义
\def \majorTitle   {综合设计报告} % 大标题
\def \minorTitleCN {基于 Tello 多旋翼无人机的智能人机交互平台} % 中文标题
\def \minorTitleEN {Intelligent Human-Computer Interaction Platform Based on Tello Multi-Rotor UAV} % 英文标题
\def \currentDate  {\zhtoday} % 日期 \zhtoday获取当前中文日期

\def \classificationNumber  {A0001} % 分类号
\def \UDC                   {} % UDC
\def \confidentialLevel     {公开} % 密级
\def \serialNumber          {0001} % 编号

% 个人信息定义
\def \titlePageInfoBox{
    % 参数：#1下划线长度 #2字号 #3标题 #4内容
    \infobox{6cm}{-2}{学~~~~~~~~院}{电子信息学院}\\
    \infobox{6cm}{-2}{专~~~~~~~~业}{通信工程}\\
    \infobox{6cm}{-2}{班~~~~~~~~级}{08062002}\\
    \infobox{6cm}{-2}{姓~~~~~~~~名}{敖冠舒}\\
    \infobox{6cm}{-2}{学~~~~~~~~号}{2020301928}\\
}

%%%%%%%%%%%%%%%%%%%%%%%%  文档开始  %%%%%%%%%%%%%%%%%%%%%%%%

\begin{document}

% 封面
\CoverPage
    {empty} % 封面类型：both、left、right、empty
    {1} % 大标题字号
    {2} % 中文标题字号
    {-2} % 英文标题字号

%%%%%%%%%%%%%%%%%%%%%  正文前页眉页脚  %%%%%%%%%%%%%%%%%%%%%%

% 页眉（关闭页眉务必将页眉类型设为empty）
\Header
    {common} % 页眉类型：common、publish、empty
    {1pt} % 上分隔线宽度
    {1pt} % 两线距离
    {0.5pt} % 下分割线宽度
    {} % 页眉左自定义内容（文本或图片）
    {\includegraphics[width=0.15\textwidth]{NWPU-title-CN.pdf}} % 页眉中自定义内容（文本或图片）
    {} % 页眉右自定义内容（文本或图片）

%============================================%

% 页脚（关闭页脚务必将页脚类型设为empty） 
\Footer
    {common} % 页脚类型：common、publish、empty
    {0pt} % 上分隔线宽度
    {0pt} % 两线距离
    {0pt} % 下分割线宽度
    {} % 页脚左自定义内容（文本或图片）
    {\thepage} % 页脚中自定义内容（文本或图片）
    {} % 页脚右自定义内容（文本或图片）

%============================================%

% 页数样式 参数：#1起始页数
\SetRomanPageNumber{1} % 设置罗马数字页码
% \setArabicPageNumber{1} % 设置阿拉伯数字页码

%%%%%%%%%%%%%%%%%%%%%%%%  摘要  %%%%%%%%%%%%%%%%%%%%%%%

\begin{abstractCN}[-2] % 中文摘要，参数：#1中文摘要标题字号

本课程设计报告标题为“基于 Tello 多旋翼无人机的智能人机交互平台”，涵盖了西北工业大学《智能人机交互与视觉感知综合设计》（U08M81002）课程的核心要求。本设计使用 Python 语言及 QtDesigner，主要在 Windows 11 环境下开发，兼具良好的移植性，可应用于其他系统。

该设计平台可以实现对大疆公司的 Tello 系列无人机的精细控制，如利用无人机的摄像头和红外等传感器实现目标检测与跟踪（如亮黄色物体、人脸、肢体等）、手势控制、体态控制等视觉功能。同时，平台还能指挥 Tello 进行拍照和一些花式动作。

Tello无人机是大疆创新与Intel公司合作开发的一款面向初学者和儿童的小型无人机，其体积小巧，重量轻，可在室内外飞行。尤其引人注目的是，其价格亲民，仅为 99 美元，且提供 Python、Java、Scratch 等多种编程语言的 SDK，成为初学者理想的无人机。

本项目以其良好的移植性，强大的可拓展性，以及低耦合的模块设计，展示了计算机、信息处理、自动控制等多学科交叉的魅力。对于将来深入研究或工业应用，此项目无疑具有非常高的价值。

% 中文关键词
\def\keywordsCN{关多旋翼无人机、计算机视觉、Python开发、自动控制}

\end{abstractCN}



%%%%%%%%%%%%%%%%%%%%%%%  启用目录  %%%%%%%%%%%%%%%%%%%%%%%%

\contentPage{1.5}{目~~~~录}            % 总目录，参数：#1目录行距 #2目录标题
\contentpageOfFigures{1.5}{图目录}     % 图目录，参数：#1目录行距 #2目录标题
\contentpageOfTables{1.5}{表目录}      % 表目录，参数：#1目录行距 #2目录标题

%%%%%%%%%%%%%%%%%%%%%%%  启用水印  %%%%%%%%%%%%%%%%%%%%%%%%

\imageWatermark % 图片水印
    {0} % 旋转角度
    {0.8} % 放缩倍率
    {0.02} % 透明度 在0-1之间调整
    {NWPU-logo.eps} % 图片路径

%%%%%%%%%%%%%%%%%%%%%  正文页眉页脚  %%%%%%%%%%%%%%%%%%%%%%

% 页眉（关闭页眉务必将页眉类型设为empty）
\Header
    {common} % 页眉类型：common、publish、empty
    {1pt} % 上分隔线宽度
    {1pt} % 两线距离
    {0.5pt} % 下分割线宽度
    {课程设计报告} % 页眉左自定义内容（文本或图片）
    {} % 页眉中自定义内容（文本或图片）
    {\currentChapterInfo} % 页眉右自定义内容（文本或图片）

%============================================%

% 页脚（关闭页脚务必将页脚类型设为empty） 
\Footer
    {common} % 页脚类型：common、publish、empty
    {0pt} % 上分隔线宽度
    {0pt} % 两线距离
    {0pt} % 下分割线宽度
    {} % 页脚左自定义内容（文本或图片）
    {\thepage} % 页脚中自定义内容（文本或图片）
    {} % 页脚右自定义内容（文本或图片）

%============================================%

% 页数样式 参数：#1起始页数
% \setRomanPageNumber{1} % 设置罗马数字页码
\SetArabicPageNumber{1} % 设置阿拉伯数字页码

%%%%%%%%%%%%%%%%%%%%%%%  正文  %%%%%%%%%%%%%%%%%%%%%%%%%%


\chapter{设计概述}

\section{设计目的}
本课程设计的首要目标是在实际操作中理解并掌握智能人机交互与视觉感知的理论知识，提高实践操作能力和工程实施能力。通过设计并实现一个基于 Tello 多旋翼无人机的智能人机交互平台，可以为后续的研究或实际应用提供技术支持。由于 Tello 无人机具有良好的开放性和扩展性，通过本课程设计，可以进一步拓展和优化该平台，实现更多功能，为未来的科研工作或者创新项目提供便利。

\section{设计内容}

\begin{itemize}
    \item 系统开发环境构建：在 Windows 11 下，利用 Python 和 QtDesigner 设计并实现系统的用户界面。同时，保证系统具备良好的移植性，以便将来在其他操作系统下使用。

    \item Tello无人机控制算法开发：设计并实现无人机的控制算法，包括基本的飞行动作（如起飞、降落、前进、后退、左转、右转等），以及一些高级的飞行动作。
    
    \item 视觉感知模块设计：利用无人机的摄像头等传感器，实现目标检测与跟踪（如亮黄色物体、人脸、肢体等）、手势控制、体态控制等视觉感知功能。
    
    \item 图像处理与识别：设计图像处理与识别算法，以便无人机能够准确识别并追踪目标。同时，利用这些图像处理和识别技术，可以对无人机拍摄的图片进行后期处理。
    
    \item 人机交互界面设计：设计人机交互界面，包括无人机控制界面、视觉感知结果显示界面等。这些界面需要友好、直观，以便用户能够轻松掌握和操作。
    
    \item 系统集成与测试：将以上所有模块集成为一个完整的系统，并进行全面的系统测试，包括单元测试、集成测试和系统测试，确保系统的稳定性和可靠性。
    
    \item 文档和报告撰写：编写系统设计文档、用户操作手册、系统维护手册等，以便其他开发者理解和维护系统，同时让用户可以快速掌握如何操作使用系统。
\end{itemize}

\section{应用平台}

\begin{table}[H]
  \centering
  \caption{硬件、软件环境一览表}
  \label{tab:soft-hardware}
  \begin{tabular}{@{}lcl@{}}
    \toprule[1.5pt]
                              & 指标     & \multicolumn{1}{c}{版本参数} \\ \midrule[0.8pt]
    \multirow{2}{*}{硬件环境} & CPU      & AMD R7-5800H               \\ \cmidrule[0.8pt](l){2-3}
                              & RAM      & 16 GB                         \\ \midrule[0.8pt]
    \multirow{2}{*}{软件环境} & 操作系统 & Windows 11 Pro 22H2 \\ \cmidrule[0.8pt](l){2-3}
                              & Python   & Python 3.8.16                 \\ \bottomrule[1.5pt]
  \end{tabular}
\end{table}

\section{开发工具}
\begin{table}[H]
  \centering
  \caption{开发工具一览表}
  \label{tab:dev-tool}
  \begin{tabular}{@{}lcl@{}}
    \toprule[1.5pt]
    工具 & 版本 & 用途 \\ \midrule[0.8pt]
    PyCharm & 2022.3 & 代码编写 \\ 
    Anaconda & 2022.11.1 & Python环境管理 \\ \bottomrule[1.5pt]
  \end{tabular}
\end{table}

\section{主要软件库}
\begin{table}[h]
\centering
\caption{主要软件库}
\renewcommand\arraystretch{1.2} % 定义表格行距
\begin{tabular}{|c|c|c|}
\hline
库名 & 版本 & 用途 \\
\hline
cvzone & 1.3.0 & 用于计算机视觉的高级函数和工具 \\
djitellopy & 2.5.0 & 大疆Tello无人机的Python接口 \\
keyboard & 0.13.5 & 键盘输入事件的处理 \\
mediapipe & 0.8.4.2 & 进行手势和肢体检测 \\
numpy & 1.24.3 & 数组处理和数值计算 \\
opencv-python & 4.8.0.74 & 计算机视觉和图像处理库 \\
pygame & 2.5.0 & 游戏开发，包括图形和声音处理 \\
PyQt5 & 5.15.9 & 用于创建图形用户界面（GUI） \\
tensorflow-intel & 2.13.0 & 优化了Intel硬件性能的TensorFlow版本 \\
\hline
\end{tabular}
\end{table}

\chapter{设计方案}
\section{总体方案}
\section{图传和目标检测方案}
图传和目标检测模块的主要工作是在无人机的飞行过程中进行实时的目标检测，需要使用 DJITelloPy 库来控制无人机，使用 OpenCV 的 DNN 模型来进行目标检测。通过实时的图像处理和键盘输入，实现了无人机的智能控制和目标追踪。

为了实现此模块的功能，开发者设计出了如下流程：
\begin{enumerate}
    \item 初始化无人机和目标检测模型：首先初始化无人机，并读取目标检测模型。无人机使用的是 DJI Tello，而目标检测模型使用的是 OpenCV 的 DNN 模型，基于 SSD MobileNet V3 Large 网络，模型训练的数据集是 COCO。
    \item 进行键盘输入控制：在无人机启动后，通过键盘进行控制，如左右、上下、前后移动，以及起飞、降落、拍照等操作。
    \item 实现目标检测：在无人机飞行的过程中，实时接收无人机的视频流并进行目标检测，通过 DNN 模型来识别和追踪目标。将识别到的目标以边框的形式标记出来，并在边框上方显示类别和置信度。
\end{enumerate}

然而在实际开发中，由于涉及到算法与硬件两部分内容，该方案存在着以下突出的技术难点：
\begin{enumerate}
    \item 实时性：无人机的图像传输和目标检测都要求有较高的实时性，这需要算法的处理速度要快。
    \item 无人机的控制：需要熟悉无人机的API，才能完成各种操作，如起飞、降落、移动等。
    \item 目标检测的准确性：目标检测的准确性直接影响了系统的整体性能，如何提高检测的准确性是一大挑战。
\end{enumerate}

经过不断地测试与调整，最终开发者摸索出一条可行方案：首先使用 DJITelloPy 库连接并控制 Tello 无人机，通过键盘输入，来发送飞行控制指令。紧接着使用 OpenCV 的 DNN 模型进行目标检测。这个模型基于 SSD MobileNet V3 Large 网络，它具有较好的准确性和较快的处理速度，适合用于实时目标检测。然后使用 cvzone.cornerRect() 函数来绘制目标的边框，使用 cv2.putText() 函数在边框上方显示类别和置信度。最后使用 cv2.imshow() 实时显示处理的图像，使用 cv2.waitKey(1) 来获取用户的输入，如选择降落或直接退出。

\section{键控方案}
键控部分，顾名思义就是通过获取键盘的输入来实现对无人机的控制，通过对应的键控指令发送到无人机，实现对无人机的智能控制。然而在实际开发中，由于考虑实时性和延迟等要求，看似简单的模块仍然存在着较大困难，主要是用传统的监控式方式读取键盘输入会占用极大的计算资源，难以实现实时性和低延迟的要求。为此，开发者抛弃了传统的“监控式”方案，改为“触发式”方案——不是每时每刻监控键盘的输入，而只有在键盘按下的时刻才进行操作。这样一来，极大地降低了对资源的占用，预期能取得良好的低延迟效果。

具体落实到方案实施，我们使用 KeyPressModule 获取键盘输入：在此基础上，通过 KeyPressModule 的 getKey() 方法来获取键盘的输入，如果检测到某个键被按下，那么就返回 True，否则返回 False。然后发送控制指令：根据键盘的输入，使用 Tello 对象的 send\_rc\_control() 方法发送对应的控制指令到无人机。例如，如果检测到 "UP" 键被按下，那么就将无人机向前移动，如果检测到 "DOWN" 键被按下，那么就将无人机向后移动。

特别值得一提的是，由于在实际操作中，“起飞”和“两个”操作具有较高的优先级，除此之外，当突发状况时也要及时停止无人机的电机。所以起飞对应的E键、降落对应的Q键，紧急停止电机的C键，仍采用“监控式”方案，这样能保证这三个操作的优先级，适当牺牲计算资源，保障飞行的安全性和可靠性。

\section{目标检测与跟踪方案}
“目标检测与追踪模块”的设计方案在整个基于 Tello 多旋翼无人机的智能人机交互平台项目中起着重要的作用。它的实现主要依赖于计算机视觉库OpenCV以及Dji Tello无人机的SDK。核心思路是利用计算机视觉技术进行实时目标检测与追踪，并通过无人机的SDK接口将检测结果转化为无人机的运动控制。

要想实现这一模块，关键要解决以下三方面问题：
\begin{enumerate}
    \item 颜色检测: 使用HSV色彩空间对目标颜色进行过滤是一项具有挑战性的工作，需要调整合适的阈值进行目标检测。
    \item 运动控制: 根据目标位置信息精确控制无人机的运动，需要精细的调试以达到较好的追踪效果。
    \item 实时性: 由于无人机需要实时接收控制指令进行运动，所以整个流程需要在较短的时间内完成，这就对代码的优化和计算资源提出了高要求。
\end{enumerate}

落实到具体方案中，我们将这一模块分为以下几个关键点，将它们在一个while循环中不断执行，直到用户终止。

\begin{itemize}
    \item 视频流处理: 获取视频流中的每一帧，进行色彩空间的映射，调整图像尺寸。
    \item 目标检测: 通过HSV色彩空间进行目标色彩的过滤，得到目标区域的二值化图像，再利用Canny算子进行边缘检测，并进行膨胀处理以便更好的找到目标的轮廓。
    \item 目标追踪: 找到轮廓后，计算轮廓的质心位置，通过质心位置判断无人机应当的运动方向。
    \item 无人机控制: 根据计算出的运动方向，控制无人机进行相应的运动。
    \item 展示: 将处理后的图像和原图像进行展示。
\end{itemize}

\begin{figure}[H]
    \centering
    \includegraphics[width=\linewidth]{pic/pic28}
    \caption{目标检测与跟踪方案实现流程图}
\end{figure}



\section{肢体与人脸追踪方案}
肢体与人脸追踪模块，核心思路是通过计算机视觉技术，捕捉到人的肢体动作或人脸，然后利用PID控制算法，根据肢体或人脸位置的变化，计算出无人机的飞行方向和速度，最后通过无人机接口发送控制信号，实现无人机的实时跟随。

解决肢体和人脸追踪问题，除了上述的实时性问题外，还有一个更关键的问题，即稳定性。稳定性要求无人机的飞行控制需要有足够的稳定性，不能因为肢体的小范围内波动就做出大的飞行调整，这需要我们引入PID算法，还需要根据实际环境以及无人机的性能特性，进行PID参数的调整。

以肢体追踪为例，我们使用OpenCV库进行视频流的获取和处理，使用cvzone库的PoseDetector模块进行肢体动作的识别和定位，使用cvzone库的PID模块进行无人机的飞行控制。在获取到肢体位置信息后，我们使用PID控制算法，对无人机的飞行方向和速度进行调整。PID控制算法的参数需要根据实际情况进行调整，以保证无人机的飞行稳定性。在整个过程中，我们还需要注意无人机的安全性问题，对于一些特殊情况，例如无人机的电量低于一定值，我们需要让无人机降落，同时也需要提供一个紧急停止的接口，以保证在出现异常情况时，可以立即停止无人机的飞行。


接下来我们同样以肢体追踪为例，其实现方案主要包括以下关键点：
\begin{enumerate}
    \item 连接Tello无人机，获取无人机的实时视频流。
    \item 将获取的视屏流进行色彩空间的转换以及大小的调整。
    \item 利用PoseDetector模块在图片上识别并标记出人的肢体动作。
    \item 获取被标记肢体的位置信息。
    \item 利用PID控制算法，根据位置信息，计算出无人机需要调整的飞行方向和速度。
    \item 通过无人机接口发送控制信号，实现无人机的实时跟随。
\end{enumerate}

\section{手势控制方案}
我们使用OpenCV库进行视频流的获取和处理，使用cvzone库的HandDetector模块进行手势的识别，使用cvzone库的FaceDetector模块进行面部的识别。

由于手部位置和姿势变化姿态过多，如果时时刻刻读取手部信息，会对无人机的控制产生极大干扰，所以在获取到手势信息后，我们首先判断手势的有效性，只有当手部位于面部的一个特定区域内时，手势才被视为有效。然后，我们根据有效的手势，执行对应的无人机控制操作，例如向上移动、向下移动、左转、右转等。

另外，我门还需要定义一个合理的手势与操作的映射关系，使得用户可以用自然和直观的方式来控制无人机。这样能提高用户的使用体验。

上述关键点可以总结为如下步骤，并将它们放在一个While循环中：
\begin{itemize}
    \item 对无人机获取的视频流进行色彩空间的转换以及大小的调整。
    \item 使用HandDetector模块识别并标记出手的位置和姿态，同时使用FaceDetector模块识别并标记出面部的位置。
    \item 根据手部和面部的相对位置，确定手势是否有效。
    \item 根据有效的手势，执行对应的无人机控制操作。
\end{itemize}

\section{肢体控制方案}
要实现肢体控制功能，同样使用OpenCV库进行视频流的获取和处理，使用cvzone库的PoseDetector模块进行姿态的识别。

在获取到姿态信息后，我们首先识别出特定的肢体动作，例如两手平举成T形状（关闭追踪），双手举起（打开追踪模式），两手在胸前交叉（准备自拍）等。然后，根据识别出的动作，执行对应的无人机控制操作，例如打开或关闭追踪模式，自动拍照等。

肢体追踪与拍照功能之前已经论述过，在此不再赘述。本部分的关键在于“如何识别出用户的不同姿态？”

要解决这个问题，cvzone库中PoseModule模块的findDistance和findAngle函数给我们提供了很大的便利。从名字可以看出，这两个函数分别用于计算两点之间的距离和三点之间的角度。在检测用户肢体时，肢体的关键点有固定的编号（如腕部是1、2，肘部是3、4，肩部是5、6），通过以上两个函数计算出肢体的长度和角度，就可以判断出用户的肢体动作，从而发送指令给无人机。


\chapter{详细设计}
\section{总体设计}
本设计使用模块化编程的方式，先分别开发不同的功能，调试完毕后形成最终的软件形态。此设计包含不同模块，实现了不同功能，核心模块具体阐述如下：

\begin{enumerate}
    \item BodyFollowing: 该模块用于实现肢体追踪。
    \item ColorObjectTracking: 该模块用于实现物体追踪，默认设定为追踪亮黄色物体，但可以进行调整。
    \item FaceFollowing: 该模块用于识别画面中的人脸并进行跟踪。
    \item HandGestureControl: 该模块用于识别画面中的手势并控制 Tello 无人机。
    \item Mapping: 该模块用于操纵无人机并画出飞行轨迹。
    \item ObjectDetection: 该模块用于在用键盘控制无人机的同时，将画面进行实时的目标检测。
    \item SelfieDrone: 该模块用于通过姿势控制无人机，并借助肢体语言实现自拍。
    \item Surveillance: 最基本的监视器模块，包含基本的飞行控制（键控模块）和监视功能。
    \item PID控制：使用工业中常用的PID控制算法，实现对无人机的精准、快速控制。
\end{enumerate}

\section{PID控制算法}
PID是以它的三种纠正算法而命名。受控变数是三种算法（比例、积分、微分）相加后的结果，即为其输出，其输入为误差值（设定值减去测量值后的结果）或是由误差值衍生的信号。另外，PID控制器可以视为是频域系统的滤波器，是一种常用、高效的控制算法，也是此项目贯穿始终的控制方法。

\begin{figure}[H]
    \centering
    \includegraphics[width=0.6\linewidth]{pic/pic1.png}
    \caption{PID算法示意图}
\end{figure}

\section{Tello无人机的基本功能}

在\href{https://dl-cdn.ryzerobotics.com/downloads/Tello/Tello_SDK_2.0_%E4%BD%BF%E7%94%A8%E8%AF%B4%E6%98%8E.pdf}{这里}给出的官方SDK使用说明中，给出了一系列常用的无人机指令，用以控制无人机。

\begin{table}[H]
    \centering
    \caption{SDK中的常用控制命令}
    \label{tab:three-line-with-note}
    \renewcommand\arraystretch{1.3} % 定义表格行距
    \setlength{\tabcolsep}{20pt} % 定义列间宽度
    \begin{threeparttable}[c]
        \begin{tabular}{ll}
            \toprule[1.5pt]
            \textbf{命令}           & \textbf{描述}\\
            \midrule[0.8pt]
            land             & 自动降落\\
            takeoff          & 自动起飞\\
            streamon         & 打开视频流\\
            streamoff        & 关闭视频流\\
            emergency        & 紧急停止电机转动\\
            up x             & 向上飞行x厘米\\
            down x           & 向下飞行x厘米\\
            left x           & 向左飞行x厘米\\
            right x          & 向右飞行x厘米\\
            forward x        & 向前飞行x厘米\\
            back x           & 向后飞行x厘米\\
            cw x             & 顺时针旋转x度\\
            ccw x            & 逆时针旋转x度\\
            \textbf{send\_rc\_control}  & \textbf{发送遥控指令}\\
            \bottomrule[1.5pt]
        \end{tabular}
    \end{threeparttable}
\end{table}

在Demo文件夹下的ImageDemo.py和FlyingDemo.py中，笔者给出了一个简单的范例帮助理解基本的控制逻辑。

\begin{lstlisting}[language=Python, caption=ImageDemo.py代码]
# 调取tello摄像头图像的基本示例
from djitellopy import tello
import cv2

me = tello.Tello() # 创建tello对象
me.connect()  # 连接tello
print(me.get_battery()) # 打印电量

me.streamoff() # 关闭视频流
me.streamon() # 打开视频流

while True:# 主循环
    img = me.get_frame_read().frame # 读取图像
    img = cv2.resize(img, (1280, 720)) # 图像大小
    img = cv2.cvtColor(img, cv2.COLOR_BGR2RGB)  #色彩空间
    cv2.imshow("Image", img) # 显示图像
    cv2.waitKey(1) # 等待1ms

cv2.destroyAllWindows()
\end{lstlisting}

\begin{lstlisting}[language=Python, caption=FlyingDemo.py代码]
# 让Tello无人机起飞的基本示例
from djitellopy import tello
import time

me = tello.Tello()
me.connect()
print(me.get_battery())

me.takeoff()
time.sleep(5) # 等待5秒钟

# 左移
start_time = time.time() # 记录当前时间
while True: # 开始无限循环
    me.send_rc_control(0, 50, 0, 0) # 发送控制命令
    if time.time() - start_time > 2: # 检查是否已经过了2秒
        break # 如果已经过了2秒，结束循环

me.land()
\end{lstlisting}

\section{键控功能}
实现键盘对无人机的控制是本项目的重要组成部分，且对于其他模块的测试和调试也有重要作用。在键盘控制模块中，笔者使用了keyboard库，该库可以实现对键盘输入事件的处理，从而实现对无人机的控制，具体代码如下：

\begin{lstlisting}[language=Python, caption=Surveillance.py代码]
# 此程序用于实现无人机的监控功能，同时兼具键盘控制功能
'''此处略去部分代码'''

def getKeyboardInput():
    lr, fb, ud, yv = 0, 0, 0, 0
    speed = 50

    if keyboard.is_pressed('left'): 
        lr = -speed
    elif keyboard.is_pressed('right'):
        lr = speed
    if keyboard.is_pressed('up'):
        fb = speed
    elif keyboard.is_pressed('down'):
        fb = -speed
    if keyboard.is_pressed('w'):
        ud = speed
    elif keyboard.is_pressed('s'):
        ud = -speed
    if keyboard.is_pressed('a'):
        yv = -speed
    elif keyboard.is_pressed('d'):
        yv = speed
    if keyboard.is_pressed('q'):
        me.land()
    if keyboard.is_pressed('e'):
        me.takeoff()
    if keyboard.is_pressed('z'):
        cv2.imwrite(f'Resources/Images/{time.time()}.jpg', img)
        time.sleep(0.3)
    return [lr, fb, ud, yv]

while True:
    vals = getKeyboardInput()
    me.send_rc_control(vals[0], vals[1], vals[2], vals[3])
    img = me.get_frame_read().frame
    img = cv2.cvtColor(img, cv2.COLOR_BGR2RGB)  # 转换色彩空间
    img = cv2.resize(img, (1080, 720))
    cv2.imshow("Image", img)
    cv2.waitKey(1)
\end{lstlisting}

可以看到，每执行一次主循环，就通过getKeyboardInput()函数获取键盘输入事件，然后根据键盘输入事件的不同，调用不同的控制命令，从而实现对无人机的控制。同时，每次循环都会读取无人机的图像，然后将图像转换为RGB色彩空间，最后显示在屏幕上。值得注意的是，进行颜色空间的转换非常必要，否则就会得到颜色失真的图像，影响后续的识别和跟踪等功能。

除此之外，当按下z键的时候，还会将当前图像保存到Resources/Images文件夹下，实现拍照功能。

\section{示踪模块}
示踪模块的主要代码如下：
\begin{lstlisting}[language=Python, caption=Mapping.py代码]
# 操控无人机的同时用红点显示出无人机的轨迹
'''此处略去导入库的代码'''

fSpeed = 117 / 10  # Forward Speed in cm/s   (15cm/s)
aSpeed = 360 / 10  # Angular Speed Degrees/s  (50d/s)
interval = 0.25
dInterval = fSpeed * interval
aInterval = aSpeed * interval

x, y = 500, 500
a = 0
yaw = 0
    
'''此处略去连接和初始化无人机的代码'''
    
points = [(0, 0), (0, 0)]

'''此处略去读取键盘输入的代码'''
def drawPoints(img, points):
    for point in points:
        cv2.circle(img, point, 5, (0, 0, 255), cv2.FILLED)
    cv2.circle(img, points[-1], 8, (0, 255, 0), cv2.FILLED)

    cv2.putText(img, f'({(points[-1][0] - 500) / 100},{(points[-1][1] - 500) / 100})m',
                (points[-1][0] + 10, points[-1][1] + 30), cv2.FONT_HERSHEY_PLAIN, 1,
                (255, 0, 255), 1)

while True:
    vals = getKeyboardInput()
    me.send_rc_control(vals[0], vals[1], vals[2], vals[3])
    img = np.zeros((1000, 1000, 3), np.uint8)
    if points[-1][0] != vals[4] or points[-1][1] != vals[5]:
        points.append((vals[4], vals[5]))
    drawPoints(img, points)
    cv2.imshow("Output", img)
    cv2.waitKey(1)
\end{lstlisting}



可以看到，示踪模块的主要功能是在操控无人机的同时用红点显示出无人机的轨迹。为了实现这一功能，首先需要定义一些参数，如前进速度、旋转速度、时间间隔、前进距离、旋转角度等。然后，定义一个points列表，用于存储无人机的轨迹点。接着，进入主循环，读取键盘输入事件，然后调用控制命令控制无人机，同时将无人机的坐标点添加到points列表中。最后，调用drawPoints()函数，将无人机的轨迹点绘制在图像上，然后显示在屏幕上。

\section{目标检测}
目标检测模块，是在键控和图传模块上衍生而来的，其主要功能是在键盘控制无人机的同时，将画面进行实时的目标检测。目标检测模块的关键问题和难点在于实现画面的低延迟，这就要求我们使用快速且准确的模型来识别，过于常用的yolo大模型准确度固然高，但会导致掉帧甚至完全不能运行，严重影响画面的实时性。主要代码如下：

\begin{lstlisting}[language=Python, caption=ObjectDetection.py代码]
# 此处省略导入库的代码

thres = 0.55 # 设置置信度阈值
nmsThres = 0.2 # 设置非极大值抑制阈值，消除冗余的边界框

classNames = []
classFile = 'coco.names'
with open(classFile, 'rt') as f:
    classNames = f.read().split('\n')

configPath = 'ssd_mobilenet_v3_large_coco.pbtxt' # 模型配置文件
weightsPath = "frozen_inference_graph.pb" # 模型权重文件

net = cv2.dnn_DetectionModel(weightsPath, configPath)
net.setInputSize(320, 320)
net.setInputScale(1.0 / 127.5)
net.setInputMean((127.5, 127.5, 127.5))
net.setInputSwapRB(True)

'''此处省略连接无人机的代码'''

# Initialize control values
lr, fb, ud, yv = 0, 0, 0, 0
speed = 50

keys_pressed = set()

'''此处省略读取键盘输入的代码'''

while True:

    vals = getKeyboardInput()
    me.send_rc_control(vals[0], vals[1], vals[2], vals[3])

    img = me.get_frame_read().frame
    img = cv2.cvtColor(img, cv2.COLOR_BGR2RGB)

    classIds, confs, bbox = net.detect(img, confThreshold=thres, nmsThreshold=nmsThres)

    try:
        for classId, conf, box in zip(classIds.flatten(), confs.flatten(), bbox):
            cvzone.cornerRect(img, box)
            cv2.putText(img, f'{classNames[classId - 1].upper()} {round(conf * 100, 2)}',
                        (box[0] + 10, box[1] + 30), cv2.FONT_HERSHEY_COMPLEX_SMALL,
                        1, (0, 255, 0), 2)
    except:
        pass

    img = cv2.resize(img, (1080, 720))
    cv2.imshow("Image", img)
    cv2.waitKey(1)
\end{lstlisting}    

通过阅读代码可以看到，为了实现目标检测这一功能，首先需要定义一些参数，如置信度阈值、非极大值抑制阈值、类别名称等。然后，读取模型配置文件和权重文件，初始化模型。接着，进入主循环，读取键盘输入事件，然后调用控制命令控制无人机，同时读取无人机的图像，然后将图像转换为RGB色彩空间，接着调用detect()函数，实现目标检测。最后，调用cornerRect()函数，将检测到的目标绘制在图像上，然后显示在屏幕上。此项目使用的权重文件仅有约13MB大小，保证了画面显示的实时性。

\section{目标跟踪}
在目标跟踪模块中，主要是对具有特征颜色的物体进行追踪，由于该功能仅仅利用RGB相机，因此需要尽量保证特征颜色与背景颜色相区分开，在本例中使用的是亮黄色。实际实现中，首先通过opencv和cvzone等库从无人机画面中提取特征颜色，并将其进行图像增强等操作。从界面中可以看到，画面被分为九个部分，根据目标所在区域不同给变量dir赋予不同的值，从而控制无人机的移动。


\section{肢体追踪与人脸追踪}
通过调用djitellopy官方库中的get\_frame\_read()方法读取Tello无人机摄像头的每帧图像，再利用cvzone的PoseModule模块对用户的肢体关键点进行检测，通过指定参数只识别上肢，将上肢框选之后确定位置，再用cvzone库中集成的PID控制模块来实现无人机的控制。在控制的过程中，还会同步绘出PID的曲线，便于观察无人机的自动控制过程。核心代码如下：

\begin{lstlisting}[language=Python, caption=BodyFollowing.py代码]
from cvzone.PoseModule import PoseDetector
detector = PoseDetector(upBody=True)

hi, wi, = 480, 640

xPID = cvzone.PID([0.22, 0, 0.1], wi // 2)
yPID = cvzone.PID([0.27, 0, 0.1], hi // 2, axis=1)
zPID = cvzone.PID([0.00016, 0, 0.000011], 150000, limit=[-20, 15])

myPlotX = cvzone.LivePlot(yLimit=[-100, 100], char='X')
myPlotY = cvzone.LivePlot(yLimit=[-100, 100], char='Y')
myPlotZ = cvzone.LivePlot(yLimit=[-100, 100], char='Z')

while True:
    img = me.get_frame_read().frame
    img = cv2.cvtColor(img, cv2.COLOR_BGR2RGB)  # 转换色彩空间
    img = cv2.resize(img, (640, 480))
    img = detector.findPose(img, draw=True)
    lmList, bboxInfo = detector.findPosition(img, draw=True)

    xVal = 0
    yVal = 0
    zVal = 0

    if bboxInfo:
        cx, cy = bboxInfo['center']
        x, y, w, h = bboxInfo['bbox']
        area = w * h

        xVal = int(xPID.update(cx))
        yVal = int(yPID.update(cy))
        zVal = int(zPID.update(area))
        imgPlotX = myPlotX.update(xVal)
        imgPlotY = myPlotY.update(yVal)
        imgPlotZ = myPlotZ.update(zVal)

        img = xPID.draw(img, [cx, cy])
        img = yPID.draw(img, [cx, cy])
        imgStacked = cvzone.stackImages([img, imgPlotX, imgPlotY, imgPlotZ], 2, 0.75)
        cv2.putText(imgStacked, str(area), (20, 50), cv2.FONT_HERSHEY_PLAIN, 3, (255, 0, 255), 3)
    else:
        imgStacked = cvzone.stackImages([img], 1, 0.75)

    me.send_rc_control(0, -zVal, -yVal, xVal)
\end{lstlisting}
人脸追踪与肢体追踪的代码与此类似，同样是借助cvzone中的人脸识别模块，再利用PID控制无人机姿态，此处不再赘述。
    
\section{手势控制与肢体控制}
手势控制与肢体控制的代码相似，对于手势控制，笔者借助mediapipe库识别手部关节，通过读取手部手指的信息对无人机发出指令实现控制。值得一提的是，为了实现手势控制，此处还加入了人脸识别功能，目的在于手的姿态太过丰富，极其容易误操作，因此只有右手举到人脸右侧区域时才会触发检测，从而极大地增强了可操作性。

肢体控制同样调用了mediapipe库，此软件库的功能及其强大，不仅能识别关节节点，还能对关节节点进行编号，利用不同编号节点的角度、距离关系，我们可以定义不同的姿态，再结合PID控制，实现多种多样的功能如自拍，极具趣味性。其中，用于判断姿态的核心代码如下所示：
\begin{lstlisting}[language=Python, caption=姿态识别核心代码]
    if detector.angleCheck(angArmL, 90) and detector.angleCheck(angArmR, 270):
    gesture = 'Tracking Mode: OFF' # 两手平举成T形状，关闭追踪
    colorG = (0, 0, 255)
    following = False
elif detector.angleCheck(angArmL, 170) and detector.angleCheck(angArmR, 180):
    gesture = 'Tracking Mode: ON' # 双手举起，打开追踪模式
    colorG = (0, 255, 0)
    following = True
elif crossDistL:
    if crossDistL < 70 and crossDistR < 70:
        gesture = "Cross" # 两手在胸前交叉，准备自拍
        snapTimer = time.time()
\end{lstlisting}

\chapter{使用说明及完成情况}
鉴于此项目功能较丰富且具有实操性，图片不足以展示其全部功能，又因视频无法在报告中体现，详情请见线下实地操作和演示。

\section{连接与初始化}
按下无人机侧边的按键，观察到指示灯在各种颜色间切换闪烁，表示此时等待配对，如图所示：

\begin{figure}[H]
    \centering
    \includegraphics[width=0.6\linewidth]{pic/pic2}
    \caption{无人机等待配对的状态}
\end{figure}

此时，打开电脑的无线网络设置，可以看到无人机已经出现在可用网络列表中，选择并连接，如图所示：
\begin{figure}
    \centering
    \includegraphics[width=0.6\linewidth]{pic/pic3}
    \caption{选择网络并连接}
\end{figure}

连接之后，观察到无人机指示灯变为黄色快闪，表示此时无人机已经连接到电脑，如图所示：

\begin{figure}[H]
    \centering
    \includegraphics[width=0.6\linewidth]{pic/pic4}
    \caption{无人机连接到电脑后的指示灯}
\end{figure}

之后运行GUI.py文件，即可打开图形用户界面，如图所示：

\begin{figure}[H]
    \centering
    \includegraphics[width=0.6\linewidth]{pic/pic5}
    \caption{图形用户界面}
\end{figure}

\section{示踪功能}

点击Mapping按钮，即可进入示踪功能，此时可以用键盘控制无人机，其控制逻辑如下：


\begin{table}[H]
    \centering
    \caption{键盘对应的常用控制操作}
    \label{tab:three-line-with-note}
    \renewcommand\arraystretch{1.3} % 定义表格行距
    \setlength{\tabcolsep}{20pt} % 定义列间宽度
    \begin{threeparttable}[c]
        \begin{tabular}{ll}
            \toprule[1.5pt]
            \textbf{按键}           & \textbf{对应操作}\\
            \midrule[0.8pt]
            W            & 上升\\
            S            & 下降\\
            A            & 逆时针旋转\\
            D            & 顺时针旋转\\
            上            & 前进\\
            下            & 后退\\
            左            & 左移\\
            右            & 右移\\
            E            & 起飞\\
            Q            & 降落\\
            Z            & 拍照\\
            C           & 强制停止\\
            \bottomrule[1.5pt]
        \end{tabular}
    \end{threeparttable}
\end{table}

在操作无人机的同时，还可以看到无人机的轨迹，如图所示：

\begin{figure}[H]
    \centering
    \includegraphics[width=0.6\linewidth]{pic/pic6}
    \caption{示踪功能的轨迹}
\end{figure}

\section{监视和键控功能}
点击Surveillance按钮，即可进入监视功能，此时可以用键盘控制无人机，其控制逻辑同上表，同时窗口中会实时显示无人机的图像，如图所示：

\begin{figure}[H]
    \centering
    \includegraphics[width=0.6\linewidth]{pic/pic7}
    \caption{监视功能所示的画面}
\end{figure}

\begin{figure}[H]
    \centering
    \includegraphics[width=0.6\linewidth]{pic/pic8}
    \caption{起飞后无人机状态}
\end{figure}

\section{目标检测功能}
点击ObjectDetection按钮，即可进入目标检测功能，此时可以用键盘控制无人机，同时窗口中会实时显示无人机的图像，且经过目标检测算法的处理，如图所示：


\begin{figure}[H]
    \centering
        \includegraphics[width=0.45\linewidth]{pic/pic9}\hfill
        \includegraphics[width=0.45\linewidth]{pic/pic10}
        \caption{目标检测功能实现效果}
        \label{fig:multi-image-04}
    \end{figure}

\section{肢体跟踪功能}
点击BodyFollowing按钮，即可打开肢体跟踪功能，按E键起飞后，尽量保持上半身在摄像头中，如图所示：

\begin{figure}[H]
    \centering
    \includegraphics[width=0.6\linewidth]{pic/pic11}
    \caption{肢体跟踪功能}
\end{figure}

与此同时，界面上会显示出经过标记的肢体关键点，以及X、Y、Z三个方向的PID曲线，便于观察无人机的控制调节过程，如图所示：

\begin{figure}[H]
    \centering
    \includegraphics[width=0.6\linewidth]{pic/pic12}
    \caption{肢体跟踪的功能界面}
\end{figure}

\section{物体跟踪}

点击ColorObjectTracking按钮，即可进入物体跟踪功能，首先会弹出两个窗口用于设定初始数值，在此我们不进行操作，选择默认值即可，如图所示：

\begin{figure}
    \centering
    \includegraphics[width=0.6\linewidth]{pic/pic13}
    \caption{物体跟踪功能的初始设定}
\end{figure}

与此同时，界面上会显示出无人机的图像及经过增强处理后的图像，并根据目标物体的位置，实时显示出对无人机的实时控制指令，实现无人机对目标物体的跟踪，如图所示：

\begin{figure}[H]
    \centering
        \includegraphics[width=0.49\linewidth]{pic/pic14}\hfill
        \includegraphics[width=0.49\linewidth]{pic/pic15}
        \caption{目标跟踪功能的界面}
        \label{fig:multi-image-04}
    \end{figure}

\begin{figure}
    \centering
    \includegraphics[width=0.6\linewidth]{pic/pic16}
    \caption{物体跟踪功能示例}
\end{figure}

\section{面部跟踪}
点击FaceFollowing按钮，即可进入面部跟踪功能，按E键起飞后，尽量保持面部在摄像头中。同时界面中会实时显示出无人机经过面部检测算法处理后的图像，并同步绘制出PID曲线，便于观察无人机的控制调节过程，如图所示：

\begin{figure}[H]
    \centering
    \includegraphics[width=0.6\linewidth]{pic/pic17}
    \caption{面部跟踪功能}
\end{figure}

\section{姿势控制自拍}
点击SelfieDrone按钮，即可进入姿势控制自拍功能，按E键起飞后，尽量保持上半身在摄像头中，当双手举起时，无人机会进入追踪模式，会实时跟随着用户的移动，如图所示：

\begin{figure}[H]
    \centering
    \includegraphics[width=0.49\linewidth]{pic/pic18}
    \caption{追踪模式打开}
\end{figure}

当双手放下时，追踪模式关闭，无人机会保持当前位置不再追踪，如图所示：

\begin{figure}[H]
    \centering
    \includegraphics[width=0.49\linewidth]{pic/pic19}
    \caption{追踪模式关闭}
\end{figure}

当双手在胸前交叉时，界面显示Ready，表示无人机已经准备好拍照，如图所示：

\begin{figure}[H]
    \centering
    \includegraphics[width=0.49\linewidth]{pic/pic20}
    \caption{无人机准备拍照}
\end{figure}

当双手从交叉状态分开时，无人机会拍摄一张照片存储在当前目录下，实现自拍，如图所示：

\begin{figure}[H]
    \centering
    \includegraphics[width=0.49\linewidth]{pic/pic21}
    \caption{来自无人机的自拍照}
\end{figure}

\section{手势控制}
手势控制功能分为地面端和无人机端两种模式。

\subsection{地面端}
对于地面端的手势控制程序，调用了本机的摄像头，通过mediapipe库检测指尖的位置进行升降操作，在操作过程中，通过界面虚拟的摇杆对无人机进行操控，如图所示：

\begin{figure}[H]
    \centering
    \includegraphics[width=0.49\linewidth]{pic/pic27}
    \caption{地面端手势控制界面}
\end{figure}


\begin{figure}[H]
    \centering
    \includegraphics[width=0.45\linewidth]{pic/pic22}\hfill
    \includegraphics[width=0.45\linewidth]{pic/pic23}
    \caption{地面端手势控制起降}
\end{figure}

\begin{figure}[H]
    \centering
    \includegraphics[width=0.49\linewidth]{pic/pic24}
    \caption{地面端手势控制方向}
\end{figure}




\subsection{无人机端}

对于无人机端的手势控制程序，调用无人机传回的实时画面，借助cvzone库识别手指的伸出与收回，从而判断用户的意图，实现对无人机的控制，特别注意的是，为了保证控制的准确性，在人脸右侧区域的手势才会被识别。具体操作如下图所示：

\begin{figure}[H]
    \centering
    \includegraphics[width=0.45\linewidth]{pic/pic25}\hfill
    \includegraphics[width=0.45\linewidth]{pic/pic26}
    \caption{无人机端手势控制界面}
\end{figure}


\chapter{设计总结}
\section{设计亮点}
\begin{enumerate}
    \item \textbf{使用无人机实物进行开发}，实现了真正的“人机交互”，程序开发并没有停留在“纸上谈兵”。
    \item 项目涉及程序开发、计算机视觉、自动控制、图像处理等多个领域，涉及面广，难度大。尤其是自动控制部分，需要对PID控制算法有深入的理解，才能实现对无人机的精准控制，极具挑战性，体现了开发者的跨学科综合能力。
    \item 项目设计功能丰富，包含了无人机的基本控制、目标检测、目标跟踪、肢体追踪、人脸追踪、手势控制、肢体控制、自拍等多种功能，且每个功能都能独立运行，具有很强的实用性和趣味性。
    \item 具有较强的可移植性和拓展性，可以在DJI的其他无人机上运行，也可以在其他平台上运行。
    \item 占用资源小，可实现在低配置的计算机上部署运行。
    \item 分布式结构，无人机仅实现采集和通信功能，所有的计算都在地面端完成，降低了无人机的成本，未来可以实现无人机编队等更多功能。
\end{enumerate}

\section{设计不足}
限于时间、开发者能力和硬件条件，本项目仍有一些不足之处，主要有以下几点：
\begin{enumerate}
    \item 无人机仅靠RGB相机来获取信息，容易受背景图像的影响。
    \item 无人机自身支持的功能仍较少，无法自行安装拓展模块，限制了更多功能的开发。
    \item 界面美观度、人性化程度、趣味性、稳定性仍有待提高。
    \item Python脚本语言运行速度仍无法与C语言相比，这限制了未来的大规模开发。
\end{enumerate}

\section{未来展望}
总体来看，本程序的功能实现较为完整，较好地实现了指定的要求，并能在基本功能上进行扩展。鉴于本项目仍有很大的拓展空间，未来可以实现的功能有：
\begin{enumerate}
    \item 加入更多的人机交互方式，如语音控制甚至脑机接口。
    \item 借助Jetson等嵌入式平台，外接深度相机、红外相机，实现更多的功能，如避障、自主导航等。
    \item 将程序移植到其他平台上，如手机、平板电脑，开发对应的APP，降低实现门槛。
    \item 探索高性能计算语言，如C++、Rust等，提高程序运行速度，挖掘大规模开发的潜力。
    \item 提高界面美观度、软件人性化程度、趣味性、稳定性等指标。
    \item 继续拓展功能，如对指定人脸或指定物体的跟踪；存储飞行路径并返航等。
    \item 将程序开源并撰写中英文文档，方便更多的开发者使用。
    \item ……
\end{enumerate}

\section{总结与感悟}
转眼间来到了大三的结束，而这门综合设计课程也陪伴了我三个学期。回想起来，从最简单的“Hello World”到代码行数上万的大工程，从对领域懵懵懂懂到有了自己的想法与见解，与这门课相伴的过程中我收获了太多。

“纸上得来终觉浅，绝知此事要躬行”，实现这个课程设计的过程让我深刻体会到“知易行难”，在今后的学习与科研中我也会更加注重实践，不断提高自己的动手能力，力求成为一名有能力、有格局的工程师。

回顾起来，这绝对是我本科阶段最用心的一门课程设计，即使有无数的困难和障碍，但当我第一次看见自己操控无人机起飞、第一次让无人机跟着我，那种喜悦之情溢于言表。我想，这就是工程师的骄傲吧。



\StartAppendix % 启用附录

\chapter{相关代码及链接}

项目源代码已经上传至\href{https://github.com/guanshuao/TelloDrone}{github.com/guanshuao/TelloDrone}，欢迎访问。

本报告使用 \LaTeX 编写，源代码已经上传至\href{https://www.overleaf.com/read/mncxpsftfkqd}{overleaf.com/read/mncxpsftfkqd}，欢迎访问。

本报告的 \LaTeX 模板也由笔者编写，模板源代码已经上传至\href{https://www.overleaf.com/read/npcyhtngxwvp}{overleaf.com/read/npcyhtngxwvp}，欢迎访问。

作者的GitHub主页：\href{https://github.com/guanshuao}{github.com/guanshuao}。

作者的个人主页：\href{https://guanshuao.github.io}{guanshuao.github.io}。

Tello无人机的开发手册请见\href{https://dl-cdn.ryzerobotics.com/downloads/Tello/Tello_SDK_2.0_%E4%BD%BF%E7%94%A8%E8%AF%B4%E6%98%8E.pdf}{https://dl-cdn.ryzerobotics.com/downloads}。


\Header
    {common} % 页眉类型：common、publish、empty
    {1pt} % 上分隔线宽度
    {1pt} % 两线距离
    {0.5pt} % 下分割线宽度
    {} % 页眉左自定义内容（文本或图片）
    {\includegraphics[width=0.15\textwidth]{NWPU-title-CN.pdf}} % 页眉中自定义内容（文本或图片）
    {} % 页眉右自定义内容（文本或图片）

%============================================%

% 页脚（关闭页脚务必将页脚类型设为empty） 
\Footer
    {common} % 页脚类型：common、publish、empty
    {0pt} % 上分隔线宽度
    {0pt} % 两线距离
    {0pt} % 下分割线宽度
    {} % 页脚左自定义内容（文本或图片）
    {\thepage} % 页脚中自定义内容（文本或图片）
    {} % 页脚右自定义内容（文本或图片）

\StartAcknowledgements % 启用致谢

感谢课程的四位授课教师：张顺、刘奇、梅少辉、戴玉超，四位老师对我的帮助不仅仅是知识层面，更重要的是他们对待学术和科研的严谨态度，成为了我的标杆和榜样。此外特别感谢梅少辉老师提供的丰富的资源和平台，让我有机会脱离传统的课堂，自由地去学习和探究。最后感谢学校和学院提供的平台和资源让我有机会接触到这门课程，让我有机会在这门课程中学习到知识，提高自己的能力。

\end{document}